% =====================================================================
%  bib.tex —— 参考文献（原书 p.299–p.310）  译者 t07
%  条目照录原书：西文题名、作者、卷期页码逐字保留；日语文献题名保留原文，
%  解说性小注译为中文。
% =====================================================================

\chapter*{参考文献}
\addcontentsline{toc}{chapter}{参考文献}
\markboth{参考文献}{}

\begin{biblist}
% [原文 p.299]
\bibentry{Abh56}{S. S. Abhyankar. On the valuations centered in a local domain. \emph{Amer. J. Math.}, 78:321--348, 1956.}
\bibentry{AdJ97}{D. Abramovich and A. J. de Jong. Smoothness, semistability and toroidal geometry. \emph{Jour. Alg. Geom.}, 6:789--802, 1997.}
\bibentry{AGZV85}{V. I. Arnold, S. M. Gusein-Zade, and A. N. Varchenko. \emph{Singularities of Differentiable Maps I--II}. Birkhäuser, 1985.}
\bibentry{AK97}{D. Abramovich and K. Karu. Weak semistable reduction in characteristic 0. preprint, 1997.}
\bibentry{Ale91}{V. A. Alexeev. General elephants on $Q$-Fano 3-folds. \emph{Comp. Math.}, 91:91--116, 1991.}
\bibentry{Art62}{M. Artin. Some numerical criteria for contractibility of curves on algebraic surfaces. \emph{Amer. J. Math.}, 84:485--496, 1962.}
\bibentry{Art66}{M. Artin. On isolated rational singularities of surfaces. \emph{Amer. J. Math.}, 88:129--136, 1966.}
\bibentry{Art69a}{M. Artin. Algebraic approximation of structures over complete local rings. \emph{Publ. Math. IHES}, 36:23--58, 1969.}
\bibentry{Art69b}{M. Artin. Algebraisation of formal moduli I. In D. C. Spencer and S. Iyanaga, editors, \emph{Global Analysis}, pages 21--72. Tokyo Univ. Press--Princeton Univ. Press, 1969.}
\bibentry{Art70}{M. Artin. Algebraisation of formal moduli II. \emph{Ann. of Math.}, 91:88--135, 1970.}
\bibentry{Art74}{M. Artin. Algebraic construction of Brieskorn resolutions. \emph{J. Alg.}, 29:330--348, 1974.}
\bibentry{Art76}{M. Artin. \emph{Deformations of Singularities}. Tata Lecture Notes, 1976.}
\bibentry{Art86}{M. Artin. Néron models. In G. Cornell and J. Silverman, editors, \emph{Arithmetic Geometry}, pages 213--230. Springer, 1986.}
% [原文 p.300]
\bibentry{AS95}{U. Angehrn and Y.-T. Siu. Effective freeness and point separation for adjoint bundles. \emph{Invent. Math.}, 122:291--308, 1995.}
\bibentry{AW97}{M. Andreatta and J. Wiśniewski. A survey on contractions of higher dimensional varieties. In \emph{Algebraic Geometry, Santa Cruz 1995}, volume 62 of \emph{Proc. Symp. Pure Math}, pages 153--183. Amer. Math. Soc., 1997.}
\bibentry{BC94}{K. Behnke and J. Christophersen. $M$-resolutions and deformations of quotient singularities. \emph{Amer. J. Math.}, 116:881--903, 1994.}
\bibentry{Ben83}{X. Benveniste. Sur l'anneau canonique de certaines variétés de dimension 3. \emph{Inv. Math.}, 73:157--164, 1983.}
\bibentry{BG71}{S. Bloch and D. Gieseker. The positivity of the Chern classes of an ample vector bundle. \emph{Inv. Math.}, 12:112--117, 1971.}
\bibentry{BK94}{M. Brion and F. Knop. Contraction and flip for varieties with group action of small complexity. \emph{J. Math. Sci. Univ. Tokyo}, 1:641--655, 1994.}
\bibentry{Bla56}{M. A. Blanchard. Sur les variétés analytiques complexes. \emph{Ann. Sci. Ec. Norm. Sup.}, 73:157--202, 1956.}
\bibentry{BP96}{F. A. Bogomolov and T. G. Pantev. Weak Hironaka theorem. \emph{Math. Res. Lett.}, 3:299--308, 1996.}
\bibentry{BPdV84}{W. Barth, C. Peters, and A. Van de Ven. \emph{Compact Complex Surfaces}. Springer, 1984.}
\bibentry{Bre97}{G. Bredon. \emph{Sheaf Theory}. Springer, 1997.}
\bibentry{Bri66}{E. Brieskorn. Über die Auflösung gewisser Singularitäten von holomorphe Abbildungen. \emph{Math. Ann.}, 166:76--102, 1966.}
% [Tutor Note: 原书题名作 von holomorphe Abbildungen，按德语语法应为 holomorphen；此系原书排印之误，照录不改。]
\bibentry{Bri68}{E. Brieskorn. Rationale Singularitäten komplexer Flächen. \emph{Inv. Math.}, 4:336--358, 1968.}
\bibentry{Bri71}{E. Brieskorn. Singular elements of semisimple algebraic groups. In \emph{Proc. Int. Congr. Math. Nice}, pages 279--284. Gauthier-Villars, 1971.}
\bibentry{BS76}{C. Bănică and O. Stănăşilă. \emph{Algebraic methods in the global theory of complex spaces}. John Wiley and Sons, 1976.}
\bibentry{CF93}{F. Campana and H. Flenner. Projective threefolds containing a smooth rational surface with ample normal bundle. \emph{J. f.r.u.a. Math. (J. Reine Angew. Math.)}, 440:77--98, 1993.}
\bibentry{CKM88}{H. Clemens, J. Kollár, and S. Mori. \emph{Higher Dimensional Complex Geometry}. Astérisque, 1988.}
% [原文 p.301]
\bibentry{Cor94}{A. Corti. Factoring birational maps of threefolds after Sarkisov. \emph{J. Alg. Geom.}, 4:223--254, 1994.}
\bibentry{Cor96}{A. Corti. Del Pezzo surfaces over Dedekind schemes. \emph{Ann. Math.}, 144:641--683, 1996.}
\bibentry{CP97}{F. Campana and Th. Peternell. Towards a Mori theory on compact Kähler threefolds. \emph{Math. Nachr.}, 187:29--59, 1997.}
\bibentry{Cut88a}{S. D. Cutkosky. Elementary contractions of Gorenstein threefolds. \emph{Math. Ann.}, 280:521--525, 1988.}
\bibentry{Cut88b}{S. D. Cutkosky. Weil divisors and symbolic algebras. \emph{Duke Math. J.}, 57:175--183, 1988.}
\bibentry{Dan82}{V. I. Danilov. The birational geometry of toric 3-folds. \emph{Izv. Akad. Nauk SSSR Ser. Mat.}, 46:972--981, 1982. (English translation) \emph{Math. USSR-Izv.}, 21:269--280, 1983.}
\bibentry{DG94}{I. Dolgachev and M. Gross. Elliptic threefolds I: Ogg--Shafarevich theory. \emph{J. Alg. Geom.}, 3:39--80, 1994.}
\bibentry{DH98}{I. Dolgachev and Y. Hu. Variation of geometric invariant theory quotient. \emph{Publ. Math. IHES}, to appear, 1998.}
\bibentry{DM69}{P. Deligne and D. Mumford. The irreducibility of space of curves of given genus. \emph{Publ. Math. IHES}, 36:75--110, 1969.}
\bibentry{Dur79}{A. Durfee. Fifteen characterizations of rational double points and simple critical points. \emph{L'Ens. Math.}, 25:131--163, 1979.}
\bibentry{DV34}{P. Du Val. On isolated singularities of surfaces which do not affect the conditions of adjunction I--III. \emph{Proc. Camb. Phil. Soc.}, 30:453--465, 483--491, 1934.}
\bibentry{EL97}{L. Ein and R. Lazarsfeld. Singularities of theta divisors and the birational geometry of irregular varieties. \emph{Jour. AMS}, 10:243--258, 1997.}
\bibentry{Elk78}{R. Elkik. Singularités rationnelles et déformations. \emph{Inv. Math.}, 47:139--147, 1978.}
\bibentry{Elk81}{R. Elkik. Rationalité des singularités canoniques. \emph{Inv. Math.}, 64:1--6, 1981.}
\bibentry{EV92}{H. Esnault and E. Viehweg. \emph{Lectures on vanishing theorems}. Birkhäuser, 1992.}
% [原文 p.302]
\bibentry{Fle81}{H. Flenner. Rational singularities. \emph{Arch. Math.}, 36:35--44, 1981.}
\bibentry{Fle89}{A. Fletcher. Working with weighted complete intersections. preprint, 1989.}
\bibentry{Fuj83}{T. Fujita. Semipositive line bundles. \emph{J. Fac. Sci. Univ. Tokyo}, 30:353--378, 1983.}
\bibentry{Fuj85}{T. Fujita. A relative version of Kawamata-Viehweg's vanishing theorem. preprint, 1985.}
\bibentry{Fuj86}{T. Fujita. Zariski decomposition and canonical rings of elliptic threefolds. \emph{J. Math. Soc. Japan}, 38:19--37, 1986.}
\bibentry{Ful84}{W. Fulton. \emph{Intersection Theory}. Springer, 1984.}
\bibentry{Fur86}{M. Furushima. Singular Del Pezzo surfaces. \emph{Nagoya Math. J.}, 104:1--28, 1986.}
\bibentry{GH78}{P. Griffiths and J. Harris. \emph{Principles of Algebraic Geometry}. John Wiley and Sons, 1978.}
\bibentry{God58}{R. Godement. \emph{Topologie algébrique et théorie des faisceaux}. Hermann, 1958.}
\bibentry{GR70}{H. Grauert and O. Riemenschneider. Verschwindungssätze für analytische Kohomologiegruppen auf komplexen Räumen. \emph{Invent. Math.}, 11:263--292, 1970.}
\bibentry{GR84}{H. Grauert and R. Remmert. \emph{Coherent analytic sheaves}. Springer, 1984.}
\bibentry{Gra62}{H. Grauert. Über Modifikationen und exzeptionelle analytische Mengen. \emph{Inv. Math.}, 146:331--368, 1962.}
% [Tutor Note: 原书印 Inv. Math. 146，实为 Math. Ann. 146：Invent. Math. 迟至 1966 年创刊，1962 年尚无此刊。此处照录原书，引用时宜核原刊。]
\bibentry{Gra72}{H. Grauert. Über die Deformationen isolierter Singularitäten analytischer Mengen. \emph{Inv. Math.}, 15:171--198, 1972.}
\bibentry{Gro68}{A. Grothendieck. \emph{Cohomologie Locale des Faisceaux Cohérents et Théorèmes de Lefschetz Locaux et Globaux --- SGA 2}. North-Holland, 1968.}
\bibentry{Har66}{R. Hartshorne. \emph{Residues and Duality}. Springer Lecture Notes vol. 20, 1966.}
\bibentry{Har70}{R. Hartshorne. \emph{Ample subvarieties of Algebraic Varieties}. Springer Lecture Notes vol. 156, 1970.}
% [原文 p.303]
\bibentry{Har77}{R. Hartshorne. \emph{Algebraic Geometry}. Springer, 1977.}
\bibentry{Hen11}{A. Henderson. \emph{The twenty-seven lines upon the cubic surface}. Cambridge Univ. Press, 1911.}
\bibentry{Hir60}{H. Hironaka. \emph{On the theory of birational blowing up}. PhD thesis, Harvard Univ., 1960.}
\bibentry{Hir64}{H. Hironaka. Resolution of singularities of an algebraic variety over a field of characteristic zero. \emph{Ann. of Math.}, 79:109--326, 1964.}
\bibentry{Hir71}{F. Hirzebruch. \emph{The Hilbert modular groups, resolution of the singularities at the cusps and related problems}, pages 275--288. Springer Lecture Notes vol. 244, 1971.}
\bibentry{HS71}{P. Hilton and U. Stammbach. \emph{A course in homological algebra}. Springer, 1971.}
\bibentry{Iit77}{S. Iitaka. On logarithmic Kodaira dimension of algebraic varieties. In \emph{Complex analysis and algebraic geometry}, pages 175--190. Cambridge Univ. Press, 1977.}
\bibentry{Isk80}{V. A. Iskovskikh. Anticanonical models of three-dimensional algebraic varieties. \emph{J. Soviet Math.}, 13:745--814, 1980.}
\bibentry{K+92}{J. Kollár et al. \emph{Flips and Abundance for Algebraic Threefolds}. Soc. Math. France, Astérisque vol. 211, 1992.}
\bibentry{Kaw82}{Y. Kawamata. A generalisation of Kodaira-Ramanujam's vanishing theorem. \emph{Math. Ann.}, 261:43--46, 1982.}
\bibentry{Kaw84a}{Y. Kawamata. The cone of curves of algebraic varieties. \emph{Annals of Math.}, 119:603--633, 1984.}
\bibentry{Kaw84b}{Y. Kawamata. On the finiteness of generators of the pluri-canonical ring for a threefold of general type. \emph{Amer. J. Math.}, 106:1503--1512, 1984.}
\bibentry{Kaw86}{Y. Kawamata. On the plurigenera of minimal algebraic threefolds with $K\approx0$. \emph{Math. Ann.}, 275:539--546, 1986.}
\bibentry{Kaw88}{Y. Kawamata. The crepant blowing-up of 3-dimensional canonical singularities and its application to the degeneration of surfaces. \emph{Annals of Math.}, 127:93--163, 1988.}
\bibentry{Kaw91}{Y. Kawamata. On the length of an extremal rational curve. \emph{Inv. Math.}, 105:609--611, 1991.}
% [原文 p.304]
\bibentry{Kaw92a}{Y. Kawamata. Boundedness of $Q$-Fano threefolds. In \emph{Proc. Int. Conf. Algebra}, pages 439--445. Contemp. Math. vol. 131, 1992.}
\bibentry{Kaw92b}{Y. Kawamata. Termination of log flips for algebraic 3-folds. \emph{Internat. J. Math.}, 3:653--660, 1992.}
\bibentry{Kaw94}{Y. Kawamata. Semistable minimal models of threefolds in positive or mixed characteristic. \emph{J. Alg. Geom.}, 3:463--491, 1994.}
\bibentry{Kaw97}{Y. Kawamata. On Fujita's freeness conjecture for 3-folds and 4-folds. \emph{Math. Ann.}, 308:491--505, 1997.}
\bibentry{Kaw98}{Y. Kawamata. Deformations of canonical singularities. to appear, 1998.}
\bibentry{Kee98}{S. Keel. Basepoint freeness for nef and big line bundles in positive characteristic. \emph{Annals of Math.}, to appear, 1998.}
\bibentry{KKMSD73}{G. Kempf, F. Knudsen, D. Mumford, and B. Saint-Donat. \emph{Toroidal embeddings I}. Springer Lecture Notes vol. 339, 1973.}
\bibentry{Kle66}{S. Kleiman. Toward a numerical theory of ampleness. \emph{Annals of Math.}, 84:293--344, 1966.}
\bibentry{KM87}{Y. Kawamata and K. Matsuki. The number of the minimal models for a 3-fold of general type is finite. \emph{Math. Ann.}, 276:595--598, 1987.}
\bibentry{KM92}{J. Kollár and S. Mori. Classification of three dimensional flips. \emph{Jour. AMS}, 5:533--703, 1992.}
\bibentry{KMM87}{Y. Kawamata, K. Matsuda, and K. Matsuki. \emph{Introduction to the Minimal Model Problem}, volume 10 of \emph{Adv. Stud. Pure Math.}, pages 283--360. Kinokuniya--North-Holland, 1987.}
\bibentry{KMM94a}{S. Keel, K. Matsuki, and J. McKernan. Log abundance theorem for threefolds. \emph{Duke Math. J.}, 75:99--119, 1994.}
\bibentry{KMM94b}{J. Kollár, Y. Miyaoka, and S. Mori. Boundedness of $Q$-Fano threefolds. unpublished, 1994.}
\bibentry{Kol83}{J. Kollár. \emph{Toward moduli of singular varieties}. PhD thesis, Brandeis Univ., 1983.}
\bibentry{Kol84}{J. Kollár. The cone theorem. \emph{Annals of Math.}, 120:1--5, 1984.}
\bibentry{Kol86a}{J. Kollár. Higher direct images of dualizing sheaves I. \emph{Annals of Math.}, 123:11--42, 1986.}
% [原文 p.305]
\bibentry{Kol86b}{J. Kollár. Higher direct images of dualizing sheaves II. \emph{Annals of Math.}, 124:171--202, 1986.}
\bibentry{Kol89}{J. Kollár. Flops. \emph{Nagoya Math. J.}, 113:15--36, 1989.}
\bibentry{Kol91a}{J. Kollár. Extremal rays on smooth threefolds. \emph{Ann. Sci. Ec. Norm. Sup.}, 24:339--361, 1991.}
\bibentry{Kol91b}{J. Kollár. Flips, flops, minimal models, etc. \emph{Surv. in Diff. Geom.}, 1:113--199, 1991.}
\bibentry{Kol92}{J. Kollár. Cone theorems and bug-eyed covers. \emph{J. Alg. Geom.}, 1:293--323, 1992.}
\bibentry{Kol94a}{J. Kollár. \emph{Log surfaces of general type; some conjectures}, volume 162 of \emph{Contemp. Math.}, pages 261--275. AMS, 1994.}
\bibentry{Kol94b}{J. Kollár. Push forward and base change for open immersions. Salt Lake City summer school on Moduli of Surfaces, 1994.}
\bibentry{Kol95}{J. Kollár. \emph{Shafarevich Maps and Automorphic Forms}. Princeton Univ. Press, 1995.}
\bibentry{Kol96}{J. Kollár. \emph{Rational Curves on Algebraic Varieties}, volume 32 of \emph{Ergebnisse der Math.} Springer, 1996.}
\bibentry{Kol97}{J. Kollár. Singularities of pairs. In \emph{Algebraic Geometry, Santa Cruz 1995}, volume 62 of \emph{Proc. Symp. Pure Math.}, pages 221--287. Amer. Math. Soc., 1997.}
\bibentry{KSB88}{J. Kollár and N. Shepherd-Barron. Threefolds and deformations of surface singularities. \emph{Inv. Math.}, 91:299--338, 1988.}
\bibentry{Lam86}{K. Lamotke. \emph{Regular solids and isolated singularities}. Vieweg, 1986.}
\bibentry{Lau73}{H. Laufer. Taut two-dimensional singularities. \emph{Math. Ann.}, 205:131--164, 1973.}
\bibentry{Lau77}{H. Laufer. On minimally elliptic singularities. \emph{Amer. J. Math.}, 99:1257--1295, 1977.}
\bibentry{Laz96}{R. Lazarsfeld. Lectures on linear series. In \emph{Complex Algebraic Geometry -- Park City}, volume 3 of \emph{Park City/IAS Math. Ser.}, pages 161--220. Amer. Math. Soc., 1996.}
\bibentry{Lip69}{J. Lipman. Rational singularities with applications to algebraic surfaces and unique factorization. \emph{Publ. Math. IHES}, 36:195--279, 1969.}
% [原文 p.306]
\bibentry{Loo84}{E. Looijenga. \emph{Isolated singular points on complete intersections}. Cambridge Univ. Press, 1984.}
\bibentry{Luo94}{T. Luo. Plurigenera of regular threefolds. \emph{Math. Zeitschr.}, 217:37--46, 1994.}
\bibentry{Mat69}{H. Matsumura. \emph{Commutative Algebra}. Benjamin/Cummings, 1969.}
\bibentry{Mat86}{H. Matsumura. \emph{Commutative Ring Theory}. Cambridge Univ. Press, 1986.}
\bibentry{Miy87}{Y. Miyaoka. Deformation of a morphism along a foliation. Algebraic geometry Bowdoin 1985. \emph{Proc. Symp. Pure Math.}, 46:245--268, 1987.}
\bibentry{MM81}{S. Mori and S. Mukai. Classification of Fano threefolds with $b_2\ge2$. \emph{Manuscr. Math.}, 36:147--162, 1981.}
\bibentry{MM86}{Y. Miyaoka and S. Mori. A numerical criterion for uniruledness. \emph{Annals of Math.}, 124:65--69, 1986.}
\bibentry{Mor79}{S. Mori. Projective manifolds with ample tangent bundles. \emph{Annals of Math.}, 110:593--606, 1979.}
\bibentry{Mor82}{S. Mori. Threefolds whose canonical bundles are not numerically effective. \emph{Annals of Math.}, 116:133--176, 1982.}
\bibentry{Mor85}{S. Mori. On 3-dimensional terminal singularities. \emph{Nagoya Math. J.}, 98:43--66, 1985.}
\bibentry{Mor87}{S. Mori. Classification of higher-dimensional varieties. In \emph{Algebraic Geometry, Bowdoin 1985}, volume 46 of \emph{Proc. Symp. Pure Math.}, pages 269--332. Amer. Math. Soc., 1987.}
\bibentry{Mor88}{S. Mori. Flip theorem and the existence of minimal models for 3-folds. \emph{Jour. AMS}, 1:117--253, 1988.}
\bibentry{Mor97}{D. Morrison. Mathematical aspects of mirror symmetry. In J. Kollár, editor, \emph{Complex Algebraic Geometry}, pages 265--340. AMS IAS/Park City Series, 1997.}
\bibentry{MP97}{Y. Miyaoka and T. Peternell. \emph{Geometry of higher dimensional algebraic varieties}. Birkhäuser, 1997.}
\bibentry{MS84}{D. Morrison and G. Stevens. Terminal quotient singularities in dimension 3 and 4. \emph{Proc. Amer. Math. Soc.}, 90:15--20, 1984.}
\bibentry{Muk95}{S. Mukai. Fano 多様体論の新展開——ベクトル束法とモジュライ問題を巡って——. 数学 47:124--144, 1995.}
% [原文 p.307]
\bibentry{Mum65}{D. Mumford. \emph{Geometric invariant theory}. Springer, 1965.}
\bibentry{MW97}{K. Matsuki and R. Wentworth. Mumford--Thaddeus principle on the moduli space of vector bundles on a surface. \emph{Intern. J. Math.}, 8:97--148, 1997.}
\bibentry{Nag60}{M. Nagata. On rational surfaces I--II. \emph{Mem. Coll. Sci. Univ. Kyoto}, 32, 33:351--370; 271--329, 1960.}
\bibentry{Nak87}{N. Nakayama. The lower semi-continuity of the plurigenera of complex varieties. In \emph{Algebraic Geometry, Sendai,1985}, volume 10 of \emph{Adv. Studies in Pure Math.}, pages 551--590. Kinokuniya--North-Holland, 1987.}
% [Tutor Note: 原书此条作 Sendai,1985，逗号后无空格，疑为排印疏漏；为存原书面貌，照录不改。]
\bibentry{Nak91}{N. Nakayama. Elliptic fibrations over surfaces. In \emph{Algebraic Geometry and Analytic Geometry}, pages 126--137. Springer, 1991.}
\bibentry{Nak98}{N. Nakayama. Invariance of the plurigenera of algebraic varieties. \emph{RIMS preprint}, 1998.}
\bibentry{NS95}{Y. Namikawa and J. Steenbrink. Global smoothing of Calabi--Yau threefolds. \emph{Inv. Math.}, 122:403--419, 1995.}
\bibentry{Par98}{K. H. Paranjape. Bogomolov--Pantev resolution---an expository account. In \emph{New trends in Algebraic Geometry, Warwick, July 1996}. Cambridge Univ. Press, 1998.}
\bibentry{Pin80}{H. Pinkham. Résolutions simultanée de points doubles rationels. In \emph{Séminaire sur les Singularités des Surfaces}, pages 179--204. Springer Lecture Notes vol. 777, 1980.}
\bibentry{PS97}{É. Picard and G. Simart. \emph{Théorie des fonctions algébriques}. Gauthier--Villars, 1897.}
\bibentry{Ree58}{D. Rees. On a problem of Zariski. \emph{Ill. J. Math.}, 2:145--149, 1958.}
\bibentry{Rei76}{M. Reid. Elliptic Gorenstein singularities of surfaces. preprint, 1976.}
\bibentry{Rei80}{M. Reid. Canonical threefolds. In A. Beauville, editor, \emph{Géométrie Algébrique Angers}, pages 273--310. Sijthoff \& Noordhoff, 1980.}
\bibentry{Rei83a}{M. Reid. Decomposition of toric morphisms. In \emph{Arithmetic and Geometry II}, volume 36 of \emph{Progress in Math.} Birkhäuser, 1983.}
\bibentry{Rei83b}{M. Reid. Minimal models of canonical threefolds. In S. Iitaka, editor, \emph{Algebraic Varieties and Analytic Varieties}, volume 1 of \emph{Adv. Stud. Pure Math.}, pages 131--180. Kinokuniya and North-Holland, 1983.}
% [原文 p.308]
\bibentry{Rei83c}{M. Reid. Projective morphisms according to Kawamata. preprint, 1983.}
\bibentry{Rei87a}{M. Reid. The moduli space of threefolds with $K=0$ may nevertheless be irreducible. \emph{Math. Ann.}, 278:329--334, 1987.}
\bibentry{Rei87b}{M. Reid. Young person's guide to canonical singularities. In \emph{Algebraic Geometry Bowdoin 1985}, volume 46 of \emph{Proc. Symp. Pure Math.} Amer. Math. Soc., 1987.}
\bibentry{SBW94}{N. I. Shepherd-Barron and P. Wilson. Singular threefolds with numerically trivial first and second Chern classes. \emph{J. Alg. Geom.}, 3:265--281, 1994.}
\bibentry{Ser56}{J.-P. Serre. Géométrie algébrique et géométrie analytique. \emph{Ann. Inst. Fourier}, 6:1--42, 1956.}
\bibentry{Sha94}{R. I. Shafarevich. \emph{Basic Algebraic Geometry}. Springer, 1994.}
\bibentry{Sho85}{V. Shokurov. The nonvanishing theorem. \emph{Izv. Akad. Nauk SSSR Ser. Mat.}, 49:635--651, 1985. (English translation) \emph{Math. USSR Izvestiya}, 26:591--604, 1986.}
\bibentry{Sho92}{V. Shokurov. 3-fold log flips. \emph{Izv. Russ. Akad. Nauk Ser. Mat.}, 56:105--203, 1992. (English translation) \emph{Russian Acad. Sci. Izv. Math.}, 40:95--202, 1993.}
\bibentry{Sie69}{C. L. Siegel. \emph{Topics in Complex Function Theory}, vols. I--III. John Wiley and Sons, 1969.}
\bibentry{Siu98}{Y.-T. Siu. Deformation invariance of plurigenera. \emph{Invent. Math.} to appear, 1998.}
\bibentry{Ste88}{J. Stevens. On canonical singularities as total spaces of deformations. \emph{Abb. Math. Sem. Univ. Hamburg}, 58:275--283, 1988.}
\bibentry{Ste93}{J. Stevens. Partial resolutions of quotient singularities. \emph{Manuscr. Math.}, 79:7--11, 1993.}
\bibentry{Sza95}{E. Szabó. Divisorial log terminal singularities. \emph{J. Math. Sci. Univ. Tokyo}, 1:631--639, 1995.}
\bibentry{Tak94}{K. Takegoshi. Higher direct images of canonical sheaves. \emph{Math. Ann.}, 303:389--416, 1994.}
\bibentry{Tau89}{K. Takeuchi. Some birational maps of Fano 3-folds. \emph{Compositio Math.}, 71:265--283, 1989.}
% [原文 p.309]
\bibentry{Tha96}{M. Thaddeus. Geometric invariant theory and flips. \emph{Jour. AMS}, 9:691--729, 1996.}
\bibentry{Tyu70}{G. Tyurina. Resolution of the singularities of flat deformations of rational double points. \emph{Funct. Anal. and Applic.}, 4:77--83, 1970.}
\bibentry{Uen75}{K. Ueno. \emph{Classification Theory of Algebraic Varieties and Compact Complex Spaces}. Springer Lecture Notes vol. 439, 1975.}
\bibentry{vdW91}{B. L. van der Waerden. \emph{Algebra}. (English translation) volume I. 7th ed., volume II. 5th ed. Springer, 1991.}
\bibentry{Vie82}{E. Viehweg. Vanishing theorems. \emph{J. f. r. u. a. Math. (J. Reine Angew. Math.)}, 335:1--8, 1982.}
\bibentry{Zar62}{O. Zariski. The theorem of Riemann-Roch for high multiples of an effective divisor on an algebraic surface. \emph{Annals of Math.}, 76:560--615, 1962.}
\end{biblist}

本书执笔之际，参考了下列日语文献。
% [Tutor Note: 本节日语文献的题名、作者、出版社与年份概依原书；条目下方缩进的小字（如代数幾何の入門書）为原书所附说明，已译为中文。]

\begin{biblist}
\bibentry{飯高 76--77}{飯高茂. 代数幾何学 I--III（岩波講座基礎数学）. 岩波書店, 1976--77.\newline 代数几何的入门书。}
\bibentry{上野 93}{上野健爾. 代数幾何（岩波講座応用数学）. 岩波書店, 1993.\newline 代数几何的入门书。}
\bibentry{永宮丸 72}{永田雅宜・宮西正宜・丸山正樹. 抽象代数幾何学（共立講座現代の数学）. 共立出版, 1972.\newline 代数几何的入门书。}
\bibentry{宮西 90}{宮西正宜. 代数幾何学（裳華房数学選書）. 裳華房, 1990.\newline 代数几何的入门书。}
\bibentry{川又 97}{川又雄二郎. 代数多様体論（共立講座 21 世紀の数学）. 共立出版, 1997.\newline 代数几何的入门书。其中有极小模型理论的简短解说。}
\bibentry{石井 97}{石井志保子. 特異点入門. シュプリンガー東京, 1997.\newline 代数几何及奇点的入门书。}
\bibentry{小田 85}{小田忠雄. 凸体と代数幾何学（紀伊國屋数学選書）. 紀伊國屋書店, 1985.\newline 环面嵌入理论的入门书。其中有环面簇情形下“森理论”的解说。}
% [原文 p.310]
\bibentry{永田 74}{永田雅宜. 可換環（紀伊國屋数学選書）. 紀伊國屋書店, 1974.\newline 环论的入门书。}
\bibentry{松村 80}{松村英之. 可換環（共立講座現代の数学）. 共立出版, 1980.\newline 环论的入门书。}

\end{biblist}
